% Figure blocks for Overleaf; insert only selected figures.
% Requires \usepackage{graphicx}. Copy the figures directory into Overleaf.

\begin{figure}[htbp]
  \centering
  \includegraphics[width=\linewidth]{figures_academic/MTE16_success_greedy_stochastic.pdf}
  \caption{Final-checkpoint success rate under greedy and stochastic action selection on MTE16. Bars represent the mean across five training seeds (42–46); dots are individual seed values. Each policy is evaluated over 100 episodes per action-selection mode.}
  \label{fig:MTE16_success_greedy_stochastic}
\end{figure}

\begin{figure}[htbp]
  \centering
  \includegraphics[width=\linewidth]{figures_academic/MTE16_reward_greedy_stochastic.pdf}
  \caption{Final-checkpoint return under greedy and stochastic action selection on MTE16. Bars represent the mean across five training seeds (42–46); dots are individual seed values. Each policy is evaluated over 100 episodes per action-selection mode.}
  \label{fig:MTE16_reward_greedy_stochastic}
\end{figure}

\begin{figure}[htbp]
  \centering
  \includegraphics[width=\linewidth]{figures_academic/MTE16_train_reward_last10pct.pdf}
  \caption{Mean episode return during the final 10\% of training on MTE16. Each dot represents one of five training seeds; bars show the across-seed mean.}
  \label{fig:MTE16_train_reward_last10pct}
\end{figure}

\begin{figure}[htbp]
  \centering
  \includegraphics[width=\linewidth]{figures_academic/MTE16_unique_states_final.pdf}
  \caption{Recorded unique-state counts at the end of training on MTE16 (logarithmic axis). Dots denote individual seeds and horizontal lines denote seed medians. This metric follows the implementation-specific logging definition and is not a measure of successful coordination.}
  \label{fig:MTE16_unique_states_final}
\end{figure}

\begin{figure}[htbp]
  \centering
  \includegraphics[width=\linewidth]{figures_academic/MTE16_stochastic_success_by_seed.pdf}
  \caption{Final-checkpoint stochastic task success on MTE16, disaggregated by training seed (42–46). Each dot represents one policy evaluated over 100 episodes. Seeds appear left to right in ascending order within each method.}
  \label{fig:MTE16_stochastic_success_by_seed}
\end{figure}

\begin{figure}[htbp]
  \centering
  \includegraphics[width=\linewidth]{figures_academic/PushBox16_success_greedy_stochastic.pdf}
  \caption{Final-checkpoint success rate under greedy and stochastic action selection on PushBox16. Bars represent the mean across five training seeds (42–46); dots are individual seed values. Each policy is evaluated over 100 episodes per action-selection mode.}
  \label{fig:PushBox16_success_greedy_stochastic}
\end{figure}

\begin{figure}[htbp]
  \centering
  \includegraphics[width=\linewidth]{figures_academic/PushBox16_reward_greedy_stochastic.pdf}
  \caption{Final-checkpoint return under greedy and stochastic action selection on PushBox16. Bars represent the mean across five training seeds (42–46); dots are individual seed values. Each policy is evaluated over 100 episodes per action-selection mode.}
  \label{fig:PushBox16_reward_greedy_stochastic}
\end{figure}

\begin{figure}[htbp]
  \centering
  \includegraphics[width=\linewidth]{figures_academic/PushBox16_train_reward_last10pct.pdf}
  \caption{Mean episode return during the final 10\% of training on PushBox16. Each dot represents one of five training seeds; bars show the across-seed mean.}
  \label{fig:PushBox16_train_reward_last10pct}
\end{figure}

\begin{figure}[htbp]
  \centering
  \includegraphics[width=\linewidth]{figures_academic/PushBox16_unique_states_final.pdf}
  \caption{Recorded unique-state counts at the end of training on PushBox16 (logarithmic axis). Dots denote individual seeds and horizontal lines denote seed medians. This metric follows the implementation-specific logging definition and is not a measure of successful coordination.}
  \label{fig:PushBox16_unique_states_final}
\end{figure}

\begin{figure}[htbp]
  \centering
  \includegraphics[width=\linewidth]{figures_academic/PushBox16_stochastic_success_by_seed.pdf}
  \caption{Final-checkpoint stochastic task success on PushBox16, disaggregated by training seed (42–46). Each dot represents one policy evaluated over 100 episodes. Seeds appear left to right in ascending order within each method.}
  \label{fig:PushBox16_stochastic_success_by_seed}
\end{figure}

\begin{figure}[htbp]
  \centering
  \includegraphics[width=\linewidth]{figures_academic/PassSparse16_success_greedy_stochastic.pdf}
  \caption{Final-checkpoint success rate under greedy and stochastic action selection on PassSparse16. Bars represent the mean across five training seeds (42–46); dots are individual seed values. Each policy is evaluated over 100 episodes per action-selection mode.}
  \label{fig:PassSparse16_success_greedy_stochastic}
\end{figure}

\begin{figure}[htbp]
  \centering
  \includegraphics[width=\linewidth]{figures_academic/PassSparse16_reward_greedy_stochastic.pdf}
  \caption{Final-checkpoint return under greedy and stochastic action selection on PassSparse16. Bars represent the mean across five training seeds (42–46); dots are individual seed values. Each policy is evaluated over 100 episodes per action-selection mode.}
  \label{fig:PassSparse16_reward_greedy_stochastic}
\end{figure}

\begin{figure}[htbp]
  \centering
  \includegraphics[width=\linewidth]{figures_academic/PassSparse16_train_reward_last10pct.pdf}
  \caption{Mean episode return during the final 10\% of training on PassSparse16. Each dot represents one of five training seeds; bars show the across-seed mean.}
  \label{fig:PassSparse16_train_reward_last10pct}
\end{figure}

\begin{figure}[htbp]
  \centering
  \includegraphics[width=\linewidth]{figures_academic/PassSparse16_unique_states_final.pdf}
  \caption{Recorded unique-state counts at the end of training on PassSparse16 (logarithmic axis). Dots denote individual seeds and horizontal lines denote seed medians. This metric follows the implementation-specific logging definition and is not a measure of successful coordination.}
  \label{fig:PassSparse16_unique_states_final}
\end{figure}

\begin{figure}[htbp]
  \centering
  \includegraphics[width=\linewidth]{figures_academic/PassSparse16_stochastic_success_by_seed.pdf}
  \caption{Final-checkpoint stochastic task success on PassSparse16, disaggregated by training seed (42–46). Each dot represents one policy evaluated over 100 episodes. Seeds appear left to right in ascending order within each method.}
  \label{fig:PassSparse16_stochastic_success_by_seed}
\end{figure}

\begin{figure}[htbp]
  \centering
  \includegraphics[width=\linewidth]{figures_academic/PushBox16_success_per_seed.pdf}
  \caption{Final-checkpoint stochastic task success on PushBox16 for three representative methods. Bars denote the across-seed mean and dots denote individual training seeds (ordered 42–46 from left to right). MAIC achieved 0/5 successful seeds, whereas MAIC + Exact KL and MADDPG each achieved 3/5.}
  \label{fig:PushBox16_success_per_seed}
\end{figure}

\begin{figure}[htbp]
  \centering
  \includegraphics[width=\linewidth]{figures_academic/MTE16_paired_seed_response.pdf}
  \caption{Seed-wise final-checkpoint success on MTE16 with and without exact KL regularization. Light gray lines connect matching training seeds under the original and modified objective; black lines connect across-seed means. The panels correspond to MAIC and MASAC under greedy and stochastic evaluation. Mean paired differences, in panel order, are +15.8, -3.6, +30.0, and +26.8 percentage points. These descriptive differences are not claims of statistical significance.}
  \label{fig:MTE16_paired_seed_response}
\end{figure}

\begin{figure}[htbp]
  \centering
  \includegraphics[width=\linewidth]{figures_academic/PushBox_PassSparse_recorded_unique_states.pdf}
  \caption{Recorded unique-state counts at the final training step in PushBox16 and PassSparse16. Dots indicate training seeds and horizontal segments are seed medians. The logarithmic scale and the implementation-specific logging definition should be considered when comparing methods.}
  \label{fig:PushBox_PassSparse_recorded_unique_states}
\end{figure}

\begin{figure}[htbp]
  \centering
  \includegraphics[width=\linewidth]{figures_academic/MTE16_MAIC_checkpoint_stability.pdf}
  \caption{Seed-wise stochastic task success for MAIC on MTE16 at evaluated checkpoints. Light gray lines represent individual training seeds across checkpoints; the dark line denotes their arithmetic mean. Each point is based on 100 evaluation episodes per seed. Only checkpoints evaluated for all five seeds are shown.}
  \label{fig:MTE16_MAIC_checkpoint_stability}
\end{figure}

\begin{figure}[htbp]
  \centering
  \includegraphics[width=\linewidth]{figures_academic/MTE16_MAIC_ExactKL_checkpoint_stability.pdf}
  \caption{Seed-wise stochastic task success for MAIC + Exact KL on MTE16 at evaluated checkpoints. Light gray lines represent individual training seeds across checkpoints; the dark line denotes their arithmetic mean. Each point is based on 100 evaluation episodes per seed. Only checkpoints evaluated for all five seeds are shown.}
  \label{fig:MTE16_MAIC_ExactKL_checkpoint_stability}
\end{figure}

\begin{figure}[htbp]
  \centering
  \includegraphics[width=\linewidth]{figures_academic/MTE16_MASAC_checkpoint_stability.pdf}
  \caption{Seed-wise stochastic task success for MASAC on MTE16 at evaluated checkpoints. Light gray lines represent individual training seeds across checkpoints; the dark line denotes their arithmetic mean. Each point is based on 100 evaluation episodes per seed. Only checkpoints evaluated for all five seeds are shown.}
  \label{fig:MTE16_MASAC_checkpoint_stability}
\end{figure}

\begin{figure}[htbp]
  \centering
  \includegraphics[width=\linewidth]{figures_academic/MTE16_MASAC_ExactKL_checkpoint_stability.pdf}
  \caption{Seed-wise stochastic task success for MASAC + Exact KL on MTE16 at evaluated checkpoints. Light gray lines represent individual training seeds across checkpoints; the dark line denotes their arithmetic mean. Each point is based on 100 evaluation episodes per seed. Only checkpoints evaluated for all five seeds are shown.}
  \label{fig:MTE16_MASAC_ExactKL_checkpoint_stability}
\end{figure}

\begin{figure}[htbp]
  \centering
  \includegraphics[width=\linewidth]{figures_academic/PushBox16_MAIC_checkpoint_stability.pdf}
  \caption{Seed-wise stochastic task success for MAIC on PushBox16 at evaluated checkpoints. Light gray lines represent individual training seeds across checkpoints; the dark line denotes their arithmetic mean. Each point is based on 100 evaluation episodes per seed. Only checkpoints evaluated for all five seeds are shown.}
  \label{fig:PushBox16_MAIC_checkpoint_stability}
\end{figure}

\begin{figure}[htbp]
  \centering
  \includegraphics[width=\linewidth]{figures_academic/PushBox16_MAIC_ExactKL_checkpoint_stability.pdf}
  \caption{Seed-wise stochastic task success for MAIC + Exact KL on PushBox16 at evaluated checkpoints. Light gray lines represent individual training seeds across checkpoints; the dark line denotes their arithmetic mean. Each point is based on 100 evaluation episodes per seed. Only checkpoints evaluated for all five seeds are shown.}
  \label{fig:PushBox16_MAIC_ExactKL_checkpoint_stability}
\end{figure}

\begin{figure}[htbp]
  \centering
  \includegraphics[width=\linewidth]{figures_academic/PushBox16_MASAC_checkpoint_stability.pdf}
  \caption{Seed-wise stochastic task success for MASAC on PushBox16 at evaluated checkpoints. Light gray lines represent individual training seeds across checkpoints; the dark line denotes their arithmetic mean. Each point is based on 100 evaluation episodes per seed. Only checkpoints evaluated for all five seeds are shown.}
  \label{fig:PushBox16_MASAC_checkpoint_stability}
\end{figure}

\begin{figure}[htbp]
  \centering
  \includegraphics[width=\linewidth]{figures_academic/PushBox16_MASAC_ExactKL_checkpoint_stability.pdf}
  \caption{Seed-wise stochastic task success for MASAC + Exact KL on PushBox16 at evaluated checkpoints. Light gray lines represent individual training seeds across checkpoints; the dark line denotes their arithmetic mean. Each point is based on 100 evaluation episodes per seed. Only checkpoints evaluated for all five seeds are shown.}
  \label{fig:PushBox16_MASAC_ExactKL_checkpoint_stability}
\end{figure}

\begin{figure}[htbp]
  \centering
  \includegraphics[width=\linewidth]{figures_academic/PassSparse16_MAIC_checkpoint_stability.pdf}
  \caption{Seed-wise stochastic task success for MAIC on PassSparse16 at evaluated checkpoints. Light gray lines represent individual training seeds across checkpoints; the dark line denotes their arithmetic mean. Each point is based on 100 evaluation episodes per seed. Only checkpoints evaluated for all five seeds are shown.}
  \label{fig:PassSparse16_MAIC_checkpoint_stability}
\end{figure}

\begin{figure}[htbp]
  \centering
  \includegraphics[width=\linewidth]{figures_academic/PassSparse16_MAIC_ExactKL_checkpoint_stability.pdf}
  \caption{Seed-wise stochastic task success for MAIC + Exact KL on PassSparse16 at evaluated checkpoints. Light gray lines represent individual training seeds across checkpoints; the dark line denotes their arithmetic mean. Each point is based on 100 evaluation episodes per seed. Only checkpoints evaluated for all five seeds are shown.}
  \label{fig:PassSparse16_MAIC_ExactKL_checkpoint_stability}
\end{figure}

\begin{figure}[htbp]
  \centering
  \includegraphics[width=\linewidth]{figures_academic/PassSparse16_MASAC_checkpoint_stability.pdf}
  \caption{Seed-wise stochastic task success for MASAC on PassSparse16 at evaluated checkpoints. Light gray lines represent individual training seeds across checkpoints; the dark line denotes their arithmetic mean. Each point is based on 100 evaluation episodes per seed. Only checkpoints evaluated for all five seeds are shown.}
  \label{fig:PassSparse16_MASAC_checkpoint_stability}
\end{figure}

\begin{figure}[htbp]
  \centering
  \includegraphics[width=\linewidth]{figures_academic/PassSparse16_MASAC_ExactKL_checkpoint_stability.pdf}
  \caption{Seed-wise stochastic task success for MASAC + Exact KL on PassSparse16 at evaluated checkpoints. Light gray lines represent individual training seeds across checkpoints; the dark line denotes their arithmetic mean. Each point is based on 100 evaluation episodes per seed. Only checkpoints evaluated for all five seeds are shown.}
  \label{fig:PassSparse16_MASAC_ExactKL_checkpoint_stability}
\end{figure}

\begin{figure}[htbp]
  \centering
  \includegraphics[width=\linewidth]{figures_academic/curves_MTE16.pdf}
  \caption{Training trajectories on MTE16 for the six algorithms over 10,000 episodes. The left panel shows the 100-episode moving mean of return; the right panel shows recorded unique-state counts on a symmetric-logarithmic axis. Lines denote means across five training seeds, and shaded regions span the 25th–75th seed percentiles. The state-count metric uses the implementation-specific logging convention.}
  \label{fig:curves_MTE16}
\end{figure}

\begin{figure}[htbp]
  \centering
  \includegraphics[width=\linewidth]{figures_academic/curves_PushBox16.pdf}
  \caption{Training trajectories on PushBox16 for the six algorithms over 10,000 episodes. The left panel shows the 100-episode moving mean of return; the right panel shows recorded unique-state counts on a symmetric-logarithmic axis. Lines denote means across five training seeds, and shaded regions span the 25th–75th seed percentiles. The state-count metric uses the implementation-specific logging convention.}
  \label{fig:curves_PushBox16}
\end{figure}

\begin{figure}[htbp]
  \centering
  \includegraphics[width=\linewidth]{figures_academic/curves_PassSparse16.pdf}
  \caption{Training trajectories on PassSparse16 for the six algorithms over 10,000 episodes. The left panel shows the 100-episode moving mean of return; the right panel shows recorded unique-state counts on a symmetric-logarithmic axis. Lines denote means across five training seeds, and shaded regions span the 25th–75th seed percentiles. The state-count metric uses the implementation-specific logging convention.}
  \label{fig:curves_PassSparse16}
\end{figure}
